\documentclass[../../../../main.tex]{subfiles}
\begin{document}

\begin{flushright}
\footnotesize\textit{【自動文字起こし・要確認】}
\end{flushright}

[解]
\begin{equation*}
\begin{cases}
f(x) = (ax+b)h(x) + cx + a^2-a \\
g(x) = (ax-b^2)h(x) + (a-1)x + c^2
\end{cases} \quad (a \neq 0)
\end{equation*}

より、
\begin{align*}
f(x)\text{が}h(x)\text{で割り切れる} &\iff \begin{cases} c=0 \\ a(a-1)=0 \end{cases} \quad \cdots \text{\textcircled{A}} \\
g(x)\text{が \quad \ \ \ \ \ \ \ \ \ \ \ \ \ \ \ \ \ } &\iff \begin{cases} a-1=0 \\ c^2=0 \end{cases}
\end{align*}

だから、\textcircled{A}の時、$a \neq 0$から$c=0 \land a=1$で、たしかに$g(x)$は$h(x)$で割り切れる。
一方、\textcircled{A}でないとき つまり $c \neq 0$ or $a \neq 1$ の時、
たしかに$g(x)$も$h(x)$では割り切れない。

よって示された。$\blacksquare$

\end{document}
